\chapter{KULLANILAN YAZILIMLAR ve YÖNTEMLER}
\label{ch:yontem}

Bu bölümde kullanılan yazılım araçları, veri setleri ve sistemin dört adımında uygulanan yöntemler
anlatılmaktadır. Her önemli teknoloji seçiminin gerekçesi projede bir \textit{Mimari Karar Kaydı}
(ADR) olarak yazılı tutulmuştur; ara rapor tarihinde 18 ADR bulunmaktadır.

\section{Yazılım Araçları ve Teknolojiler}

Her iş için gereksinimi karşılayan en basit araç seçilmiştir. Tablo~\ref{tab:teknoloji} kullanılan
teknolojileri ve seçilme nedenlerini özetler.

\begin{table}[!ht]
\centering
\caption{Kullanılan teknolojiler ve seçilme nedenleri}
\label{tab:teknoloji}
\footnotesize
\renewcommand{\arraystretch}{1.2}
\begin{tabularx}{\textwidth}{@{}p{2.4cm} p{4.2cm} X@{}}
\toprule
\textbf{Katman} & \textbf{Teknoloji} & \textbf{Neden} \\
\midrule
Yapay zekâ & Google Gemini Flash-Lite (ücretsiz katman), \texttt{google-genai} & İyi Türkçe, JSON şemasına zorlanmış çıktı, ücretsiz \\
Gömme (embedding) & \texttt{gemini-embedding-2}, 768 boyut & Çok dilli: TR soru, TR/EN belge \\
Arka uç & Python 3.12, FastAPI, Pydantic & Tipli, hızlı; OpenAPI şemasını kendisi üretir \\
Veritabanı & PostgreSQL 18 + pgvector, psycopg 3 & Talepler, geçmiş ve vektörler tek veritabanında \\
Ön yüz & React 19, TypeScript, Vite, Tailwind CSS, shadcn/ui, TanStack Router/Query & Tipli, OpenAPI'den üretilen API istemcisi \\
Kalite & pytest, Playwright, ruff, mypy (strict), CodeQL & Her kural ve hata düzeltmesi için test \\
Teslim & Docker, GitHub Actions, Render, Neon & Tek komutla kurulum; CI geçmeden yayın yok \\
\bottomrule
\end{tabularx}
\end{table}

\section{Veri Setleri}
\label{sec:veri}

BŞEÜ'nün gerçek öğrenci talepleri kişisel veri içerdiğinden kullanılmamış; bunun yerine BŞEÜ'nün gerçek
sistemleri (OBS, SOFRA, UZEM, eduroam, Akıllı Kart) ve öğrencilerin sık karşılaştığı sorunlar üzerine
yapılan bir araştırmadan yola çıkılarak \textbf{160 talep} yazılmış ve hepsi tek tek elle
etiketlenmiştir: 80 talepli geliştirme seti (istemler bu sette geliştirilir), 40 talepli kilitli
test seti ve dağınık, yazım hatalı, iki sorunlu taleplerden oluşan 40 talepli kilitli gerçekçi set.
Her set yarı yarıya Türkçe ve İngilizcedir.


Kilitli setler hiçbir zaman düzenlenmez ve istem geliştirmede kullanılmaz; bir koruma kancası (hook)
bu dosyalardaki her değişikliği engeller. Geliştirme setinde yedi kategori dengelidir (10--12 talep);
21 talep insana yönlendirilmeyi (escalate) gerektirir.

\section{Talep Sınıflandırma Yöntemi}

Etiket kural kitabı (\texttt{docs/labels.md}) hem elle etiketlemede hem de modelin talimatında
kullanılan tek kaynaktır. Kategoriler ve sorumlu birimler Tablo~\ref{tab:kategori}'da verilmiştir.
Kategoriden birime eşleme model tarafından değil, kodda sabit bir tabloyla yapılır.

\begin{table}[!ht]
\centering
\caption{Kategoriler ve sorumlu BŞEÜ birimleri}
\label{tab:kategori}
\small
\begin{tabularx}{\textwidth}{@{}l l X@{}}
\toprule
\textbf{Kimlik} & \textbf{Kategori} & \textbf{Birim} \\
\midrule
\texttt{account\_access} & Hesap Erişimi & Bilgi İşlem Daire Başkanlığı \\
\texttt{it\_support} & Teknik Destek & Bilgi İşlem Daire Başkanlığı \\
\texttt{registration} & Ders Kaydı & Öğrenci İşleri Daire Başkanlığı \\
\texttt{academic\_records} & Akademik Kayıtlar & Öğrenci İşleri Daire Başkanlığı \\
\texttt{billing} & Harç ve Ödemeler & Öğrenci İşleri, İstatistik Disiplin ve Harçlar Şube Md. \\
\texttt{campus\_life} & Kampüs Yaşamı & Sağlık, Kültür ve Spor Daire Başkanlığı \\
\texttt{other} & Diğer & Öğrenci İşleri nöbetçi personeli \\
\bottomrule
\end{tabularx}
\end{table}

\textbf{Escalation kuralları:} güvenlik tehdidi; taciz veya kriz belirtisi; ele geçirilmiş olabilecek
hesap; kural istisnası talebi; para veya resmî kayıt hatası; not ve sınav itirazı; hukuki tehdit;
disiplin süreci. Öğrencinin yalnızca sinirli olması tek başına bir neden değildir. İki sorun içeren
bir talepte kategori öğrenciyi en acil engelleyen sorundan seçilir, escalation ise bütün sorunlara bakar.

Modelin çıktısı katı bir Pydantic şemasıyla sınırlandırılmıştır; şemaya uymayan JSON reddedilir.
Şemada \texttt{reason} alanı ilk sıradadır: model önce gerekçesini yazar, sonra etiketleri verir.

\section{Bilgi Tabanı ve Belge Arama}

Bilgi tabanı, BŞEÜ Ön Lisans ve Lisans Eğitim-Öğretim Yönetmeliği, Öğrenci İşleri ve SKS sıkça
sorulan sorular sayfaları, Bilgi İşlem duyuruları gibi resmî kaynaklardan elle derlenmiş 14 konudan
oluşur. Her konu Türkçe ve İngilizce iki belge olarak yazılmış, her belge başlık düzeyinde parçalara
(chunk) bölünmüştür. Her parça kaynağını (belge adı, resmî kaynak satırı, bağlantı) taşır.

Her parça \texttt{gemini-embedding-2} modeliyle 768 boyutlu bir vektöre dönüştürülür. Talep de aynı
şekilde vektöre çevrilir ve talep vektörü $\mathbf{q}$ ile parça vektörü $\mathbf{d}$ arasındaki
kosinüs benzerliği Eşitlik~\ref{eq:kosinus} ile hesaplanır.

\begin{equation}
  s(\mathbf{q},\mathbf{d}) = \frac{\mathbf{q}\cdot\mathbf{d}}{\lVert\mathbf{q}\rVert\,\lVert\mathbf{d}\rVert}
  = 1 - \bigl(\mathbf{q} \Leftrightarrow \mathbf{d}\bigr)
  \label{eq:kosinus}
\end{equation}

Burada $\Leftrightarrow$, pgvector'ün kosinüs uzaklığı işlecidir (\texttt{<=>}). Arama, talebin
dilindeki parçalar arasından en yakın $k = 3$ parçayı veritabanının içinde tek bir SQL sorgusuyla
bulur. 190 parça için yaklaşık arama indeksine gerek yoktur; tam tarama hızlıdır ve hiçbir parçayı
kaçırmaz.

\section{Kaynaklı Cevap Yazma ve Alıntı Denetimi}

Cevap adımında modele talep (çitlenmiş olarak) ve bulunan üç parça verilir. Model sırayla şu alanları
doldurur: \texttt{evidence} (kuraldan kelimesi kelimesine alıntı), \texttt{covered} (belgeler ana
soruyu cevaplıyor mu), \texttt{reply} (cevap), \texttt{sources} (kullanılan parça kimlikleri) ve
\texttt{general\_help} (belgelerin cevaplamadığı kısım için resmî olmadığı etiketlenen genel tavsiye).

Modelin alıntısına güvenilmez; alıntı kodla denetlenir (\texttt{domain/grounding.py}). Alıntı
cümlelere bölünür, her cümle Unicode, tırnak ve büyük/küçük harf farkları giderilerek normalleştirilir
ve bulunan belgelerin metninde aranır; tek bir cümle bile bulunamazsa alıntı geçersizdir. Bu denetim
hem modelin uydurduğu kuralları hem de öğrencinin talebe yerleştirdiği sahte ``resmî kuralları''
(prompt injection) yakalar.

\section{Otomatik Gönderim Kararı}
\label{sec:karar}

Karar, \texttt{domain/decision.py} içinde yedi açık kontrolden oluşan sade bir Python fonksiyonudur.
Cevap yalnızca kontrollerin \textbf{tamamı} geçerse otomatik gönderilir; kalan her kontrol bir neden
kodu olarak kaydedilir ve hem personel panelinde hem de öğrencinin ``Nasıl karar verildi?'' izinde
gösterilir (Şekil~\ref{fig:karar}).

\begin{figure}[!ht]
\centering
\begin{tikzpicture}[node distance=3.2mm, every node/.style={font=\scriptsize}]
  \node[kutu, minimum width=30mm] (in) {Etiketler + en iyi skor $s$\\+ taslak + getirilen parçalar};
  \node[karar, below=of in] (k1) {kategori $\neq$ diğer?};
  \node[karar, below=of k1] (k2) {escalate yok?};
  \node[karar, below=of k2] (k3) {öncelik $\neq$ acil?};
  \node[karar, below=of k3] (k4) {$s \geq 0{,}65$?};
  \node[karar, below=of k4] (k5) {covered?};
  \node[karar, below=of k5] (k6) {kaynak getirildi?};
  \node[karar, below=of k6] (k7) {alıntı birebir?};
  \node[kodkutu, below=5mm of k7, minimum width=34mm] (s) {\textbf{Otomatik gönder}};
  \coordinate (b1) at ($(k4.east)+(24mm,0)$);
  \coordinate (b2) at ($(k4.east)+(72mm,0)$);
  \node[fit=(k1.north -| b1) (k7.south -| b2), draw=s2, fill=s2!8, rounded corners=3pt, inner sep=0pt] (h) {};
  \node[align=center, font=\small, text width=44mm] at (h.center) {\textbf{Birime yönlendir}\\[4pt] neden kodları kaydedilir,\\ personel taslağı düzenleyip\\ gönderir\\[6pt]
        \scriptsize escalation yoksa ve kaynak\\\scriptsize doğrulandıysa öğrenciye\\\scriptsize ``ön bilgi'' gösterilir};
  \foreach \a/\b in {in/k1,k1/k2,k2/k3,k3/k4,k4/k5,k5/k6,k6/k7,k7/s}
    \draw[ok] (\a) -- (\b);
  \foreach \k in {k1,k2,k3,k4,k5,k6,k7}
    \draw[ok, draw=s2] (\k.east) -- node[etiket, above, pos=0.3] {hayır} (h.west |- \k.east);
\end{tikzpicture}
\caption{Karar akışı: her ``hayır'' bir neden kodu ekler; tek bir neden bile yönlendirme için yeterlidir}
\label{fig:karar}
\end{figure}

Benzerlik eşiği (0,65) tek başına güvenilir bir ayırıcı değildir (Bölüm~\ref{sec:arama-bulgu}); bu
nedenle yalnızca bir alt sınırdır ve asıl koruma ``kapsıyor mu'' ve ``birebir alıntı'' kontrolleridir.
Talep personele gitse bile, escalation yoksa ve taslağın kaynaklı kısmı doğrulandıysa bu kısım
öğrenciye hemen ``ön bilgi'' olarak gösterilir.

\section{Karşılaştırılan Yöntemler}

Bölümün yapay zekâ projeleri için öngördüğü gibi, aynı veri üzerinde en az iki yöntem
karşılaştırılmıştır: (1) aynı modelin düşük (\textit{LOW}) ve yüksek (\textit{HIGH}) düşünme
düzeyleri, (2) \texttt{gemini-3.5-flash-lite} ve \texttt{gemini-3.1-flash-lite} modelleri ve (3)
kararın kodla verilmesi ile yalnızca skora dayalı bir eşik. Karma arama (anahtar kelime + vektör, RRF)
kalan dönem için planlanmıştır.
